\label{key}\RequirePackage{etex}

% Vi benytter 2 sidet, a4 størrelse og dansk sprog. 
\documentclass[11pt,a4paper,twoside,danish]{memoir}


\usepackage[paper=A4,pagesize]{typearea}
   
\usepackage[utf8]{inputenc} % danske tegn
\usepackage[danish]{babel} % dansk sprogpakke

% Hvis engelsk sprog ønskes, uncomment denne
% \usepackage[english]{babel}

%%% % % % % % % PAKKER % % % % % % % % % % % % % % % % % % % % % % % 
     % % % % %
     %Pakker til fonts, farver, referencer og ligenende er placeret relevant sted i dokumentet.
     % Her findes generelle pakker som bruges overordnet i dokumentet % 
     \usepackage{footnote} % Package used to controle footnote(specially in tables etc.)
     \usepackage{amsmath}
     \usepackage{slantsc}
    
        %  \usepackage{float} % http://en.wikibooks.org/wiki/LaTeX/Floats,_Figures_and_Captions
                 \usepackage{multirow} %sammenfletning af r�kker
                % \usepackage{subfig}
                 \usepackage{longtable} % tabeller der g�r over flere r
                  \usepackage{geometry}
         \usepackage[section]{placeins} %Gør det muligt at bruge \FloatBarrier
   
    
    \usepackage[compatibility=false]{caption}
    \captionsetup{font=small,labelfont=bf,textfont=it}
    \usepackage{graphicx}
    \usepackage{afterpage}
    
    \usepackage[labelsep=quad,indention=10pt]{subfig}
    \captionsetup*[subfigure]{position=bottom}
    
    
    

    \usepackage[draft]{fixme} %hvor 'options' bl.a. inkluderer final, draft, danish, inline og margin
    
   %\usepackage[pdftex,plainpages=false,pdfpagelabels,colorlinks=true,citecolor=black,linkcolor=black,breaklinks]{hyperref}
   
  
 \usepackage{rotating}
   
  
   \usepackage{eso-pic} % til inds�ttelse af forside.
   % \usepackage{clock} % Til at lave et sejt analogt ur 

   \usepackage{rotating} % rotate stuff(EVERYTHING!11)
   %\usepackage{lipsum}  % Use \lipsum to generate Lorem Ipsum text.
      
   
   
% % % % % % % % % % % % % % % MATH STUFF % % % % % % % % % % % % % %
    %\usepackage{algorithmic}

    %\usepackage{harpoon} %Tilf�jer harpun-pile over st�rrelser i math, m.m.
 
    %\usepackage{pst-plot} % For axes   
    %\usepackage{breqn}
    %\usepackage{fourier}  %................... Roman+math - Utopia
    
  
    %\usepackage{nomentbl} % http://ctan.mackichan.com/macros/latex/contrib/nomentbl/nomentbl.pdf

    %\usepackage{psfrag} % insæt tag, mest brugt i forhold til grafer

   
   

%%% DIVERSE CUSTOMS PAKKET NED I FILER FOR OVERBLIK
    %


%Dette hax g�r s�ledes at man kan lave en itemize inden i en tabel.

\usepackage{mdwlist}
\makeatletter
\def\noVSpace{\@minipagetrue}
\newenvironment{tabItemize}{%
  \@minipagetrue%
  \begin{itemize*}%
}{\vspace{-\normalbaselineskip}%
  \end{itemize*}}
\makeatother

\newcolumntype{d}[1]{D{.}{.}{#1}}


%___________SEMI-t�ttere itemize
%eks:
%\begin{itemize} \semitightlist
%   \item
%\end{itemize}


\newcommand{\semitightlist}{%
\setlength{\itemsep}{0.2cm} \setlength{\parskip}{1pt}}


%Inds�ttelse af NAnote
\usepackage{ifthen}
\newcommand{\clearemptydoublepage}{\newpage{\pagestyle{empty}%
  \cleardoublepage{}}}
\newcommand{\NAfinalversion}{false}
\newcommand{\NAcompileimages}{true}
%% IN TEXT COMMENTING COMMANDS
%% To remove comments redefine \NAfinalversion to true in NA_RepPreamble
\ifthenelse{\equal{\NAfinalversion}{false}}
  { \newcommand\Note[1]{\textcolor{red}{\large\texttt{(NOTE: #1)}}} }
  { \newcommand\Note[1]{} }

%%%% Nomenklatur under formler
%\newcommand\nomenklatur[1]{%
%\begin{table}[htb]
%\centering
%\begin{tabular}{p{0.80\textwidth}}
%\begin{center}\small \textit{#1}\end{center}
%\end{tabular}
%\end{table}}


%___Til at vise programkode
\usepackage{listings}
\lstset{extendedchars=true, basicstyle=\ttfamily,keywordstyle=\normalfont\ttfamily, columns=flexible, numbers=left, numberstyle=\tiny, breaklines, breakatwhitespace=true, language=Matlab, morecomment=[l][\color{blue}]{\%}}

\renewcommand\ttdefault{txtt}

\newcommand\uu[1]{\underline{\underline{#1}}}
\newcommand\RA{\qquad \Rightarrow}
\newcommand\qqquad{\qquad \qquad}
\newcommand\qqqquad{\qqquad \qqquad}
\newcommand\qqqqquad{\qqqquad \qqqquad}
\newcommand\degree{^{\circ}}
\newcommand\f{\\[0.3cm]}
\renewcommand{\Re}{\operatorname{Re}}
\renewcommand{\Im}{\operatorname{Im}}
\newcommand\units[1]{\left[{#1}\right]}

%%%___Morten Kjelds hax
%Inds�ttelse af fede gr�ske vektorer
\newcommand{\gve}[1]
{\textrm{\boldmath ${#1}$ \unboldmath}\!}
%Inds�ttelse af vektor
\newcommand{\ve}[1]
{\mathbf{#1}}
%Inds�ttelse af vektor med streg over
\newcommand{\ove}[1]
{\overline\mathbf{#1}}
%Inds�ttelse af �, � og � i formeludtryk
\newcommand{\�}
{\textit{\textrm{�}}}
\newcommand{\�}
{\textit{\textrm{�}}}
\newcommand{\�}
{\textit{\textrm{�}}}

\definecolor{shadecolor}{gray}{.91}


    %
%%%%%%%%%%%%%%%%%%%%%%%%%%%%%%%%%%%%%%%%%%%%%%%%%%%%%%%%%%%%%%%%%%%%%%%%%%%%%
%                              Custom functions                             %
%                         For general use in projects                      %
%%%%%%%%%%%%%%%%%%%%%%%%%%%%%%%%%%%%%%%%%%%%%%%%%%%%%%%%%%%%%%%%%%%%%%%%%%%%%

% For use in definition of new commands:
\usepackage{xspace}           % \newcommand{\cmd}{text\xspace} ensures the correct spacing between text and other words/characters.

% For cellheight adjustments in nomenklatur-function
\usepackage{cellspace}
    \addtolength\cellspacetoplimit{0pt}
    \addtolength\cellspacebottomlimit{0pt}


%
%% Custom "paragraph" function for making line-seperated paragraphs (e.g. double newline)
%

% Make double line break with no line indent (for changing the way paragraphs are made)
\newcommand{\newpar}{$\,$ \\ $\,$ \\}
% Alias of \newpar
\newcommand{\newparagraph}{\newpar{}}

\newcommand{\subsubsubsection}[1]{%
\newpar\paragraph{#1}$\,$\\%
}


%
%% Function for creating horizontal lines between points in a recipe-like list
%
\newcommand{\recipeline}{%
\vspace{2pt}
\hrule
\vspace{2pt}
}
%
%% Function for creating a boxed, centered equation  behaving like the equation environment
%
\newcommand{\boxedeq}[2]{%
	\[\fbox{%
			\addtolength{\linewidth}{-2\fboxsep}%
			\addtolength{\linewidth}{-2\fboxrule}%
			\begin{minipage}{#2}%
				\bigskip%
				\[#1\]%
				\hspace{1cm}%
			\end{minipage}%
	}\]%
			\addtolength{\linewidth}{-2\fboxsep}%
			}
			%
			%% Function for creating a boxed, centered equation  behaving like the equation environment
			%
			\newcommand{\boxedeqlabel}[3]{%
				\[\fbox{%
			\addtolength{\linewidth}{-2\fboxrule}%
			\begin{minipage}{#2}%
				\bigskip%
				\begin{equation}#1\end{equation}%
				\label{#3}%
				\hspace{1cm}%
			\end{minipage}%
	}\]%
}

%
%% Environment for displaying a 0th, 1st or 2nd order tensor in matrix format
%
\newenvironment{tensor}
{
	\begin{pmatrix}
}
{
    \end{pmatrix}
}

%
%% For displaying vectors
%
\newenvironment{vect}
{
	\begin{Bmatrix}
}
{
    \end{Bmatrix}
}

% VECTOR IS DEFINED
\newcommand{\vectq}[1]{\ensuremath{%
%	\{#1\}
\text{\overrightharp{$#1$}}
}}

% UNIT VECTOR IS DEFINED
\newcommand{\unitvecq}[1]{\ensuremath{%
\widehat{#1}
}}

% TIME INVARIANT IS DEFINED
\newcommand{\tInvar}[1]{
\widetilde{#1}
}
%
%% For displaying matrices
%
\newenvironment{matr}
{
	\begin{bmatrix}
}
{
    \end{bmatrix}
}
\newcommand{\matrq}[1]{\ensuremath{%
%	[#1]
\underline{#1}
}}
\renewcommand{\Re}[1]{
%\text{Re}\Big\langle#1\Big\rangle%
\text{Re}\left\{#1\right\}
}
\renewcommand{\Im}[1]{
%\text{Im}\Big\langle#1\Big\rangle%
\text{Im}\left\{#1\right\}
}
%
%% Function for creating small introductions for every new chapter (italic text)
%
\newcommand{\chapintro}[1]{%
    \textit{#1} \\ \linebreak%
    \noindent%
}

%
%% Function for creating small introductions for every new section (italic text)
%
\newcommand{\secintro}[1]{%
    \textit{#1} \\ \linebreak%
    \noindent%
}

%
%% Function for creating isolated remarks (indented italic text)
%
\newcommand{\remark}[1]{%
	\begin{center}%
		\parbox{0.8\textwidth}{\textit{#1}}%
	\end{center}%
    \noindent%
}

%
%% Function for creating isolated remarks (indented italic text)
%
\newcommand{\makequote}[2]{%
	\begin{center}%
		\parbox{0.8\textwidth}{"\textit{#1}" #2}%
	\end{center}
    \noindent%
}

%%% Function for using \verytightlist - not very pretty!
%\newcommand{\verytightlist}{%
%    \setlength{\itemsep}{0.05cm} \setlength{\parskip}{1pt}}
%
%
%%% Function for almost using \tightlist without using the memoir-class
%\newcommand{\allmosttightlist}{%
%    \setlength{\itemsep}{0.1cm} \setlength{\parskip}{1pt}}
%
%
%%% Function for using \semitightlist
%\newcommand{\semitightlist}{%
%    \setlength{\itemsep}{0.2cm} \setlength{\parskip}{1pt}}
%

%
%% Comments (shown in document outer margin)
%
\renewcommand{\comment}[2]{%
    \textcolor{blue}{\textbf{*}} \marginpar{%
        \fcolorbox{white}{blue}{\parbox{3.5cm}{%
            \color{white} \small \sffamily \textbf{#1:} #2}%
        }%
    }%
}
%\newcommand{\comment}[2]{} % Use for hiding comments

%
%% Fixmes (shown in document outer margin)
%
\renewcommand{\fixme}[2]{%
    \textcolor{red}{\textbf{!}} \marginpar{%
        \fcolorbox{white}{red}{\parbox{3.5cm}{%
            \color{white} \small \sffamily \textbf{FIXME:} #2 (#1)}%
        }%
    }%
}
%\newcommand{\fixme}[2]{} % Use for hiding fixmes


%
%% Pretext style headers/footers (for use before \mainmatter}
%% NB: Use only for report/book style in documentclass
%
\newenvironment{pretext}
{
	% Don't use
	\pagestyle{plain}
}
{
	\cleardoublepage
	\pagestyle{fancy}
}


%%% Function for using \bf to write in bold
\renewcommand{\bf}{%
    \bfseries%
}

%%% Function for using \it to write in italic
\renewcommand{\it}{%
    \itshape%
}


%% Home made commands by Rene
\newcommand{\abs}[1]{\lvert#1\rvert}
\newcommand{\sgn}[1]{\text{sgn}\left(#1\right)} %Sign function
\newcommand{\nrcs}[1]{\: \put(5.5,3.5){\circle{14}}#1 \,} %Laver en cirkel omkring et to cifferet tal, nummer cirkel stor fx \nrcs{15}
\newcommand{\nrc}[1]{\: \put(3,3.5){\circle{12}}#1 \,} %Laver en cirkel omkring et �t cifferet tal, nummer cirkel fx \nrc{5}
\newcolumntype{R}[1]{>{\raggedleft\arraybackslash}p{#1}} %R{2cm} laver en right orienteret s�jle med en bredde p� 2cm






%%%%%%%%%%%%%%%%%%%%%%%%%%%%%%%%%%%%%%%%%%%%%%%%%%%%%%%%%%%%%%%%%%%%%%%%%%%%%
%                              Custom functions                             %
%                       Primarily for this project; I6                      %
%%%%%%%%%%%%%%%%%%%%%%%%%%%%%%%%%%%%%%%%%%%%%%%%%%%%%%%%%%%%%%%%%%%%%%%%%%%%%

%
%% Adds vertical space around \hline in tabular, T = top
%
\newcommand\T{%
    \rule{5pt}{5ex}%
}

%
%% Adds vertical space around \hline in tabular, B = buttom
%
\newcommand\B{%
    \rule{0pt}{2.6ex}%
}

%
%% Nomenclature below equations
%
\newcommand\nomenklaturCellHeight{%
\includegraphics[scale=1]{styles/nomenklaturCellHeight.PNG}%
}
\renewcommand\nomenklaturCellHeight{}


\newcommand\nomenklatur[2]{%
\nomenclature{$#1$}{#2.}{}{}%
\begin{minipage}[top]{0.7\textwidth}\flushleft%\renewcommand\arraystretch{1.25}% (Value=1.0 is for standard spacing}%
\begin{tabular}{S{p{1.0cm}}S{p{1.3cm}} Sl}%
{$\,$} & \ensuremath{#1} &\nomenklaturCellHeight\parbox{12cm}{\footnotesize{#2.}}%
\end{tabular}%
\end{minipage}%
%}
\\ \noindent%
}

%
%% Start nomenclature below equations - use this command for the first line of each nomenclature-block.
%
\newcommand\nomenklaturstart[2]{%
\nomenclature{$#1$}{#2.}{}{}%
\begin{minipage}[top]{0.7\textwidth}\flushleft%\renewcommand\arraystretch{1.25}% (Value=1.0 is for standard spacing}%
\begin{tabular}{>{\setlength{\parindent}{-0.2cm}}S{p{1.0cm}}S{p{1.3cm}} Sl}
where:\nomenklaturCellHeight & \ensuremath{#1} &\nomenklaturCellHeight\parbox{12cm}{\footnotesize{#2.}}%
\end{tabular}%
\end{minipage}%
%}
\\ \noindent%
}

%%%%%%%%%%%%%%%%%%%%%%%%%%%%%%%%%%%%%%%%%%%%%%%%%%%%%%%%%%%
%%figur nomenklatur
%%%%%%%%%%%%
%\newcommand\ehnomenklatur[3]{%
%\noindent
%\nomenclature{$#1$}{#2}{}{}%
%\begin{minipage}[top]{0.7\textwidth}\flushleft%\renewcommand\arraystretch{1.25}% (Value=1.0 is for standard spacing}%
%\begin{tabular}{S{p{1.0cm}} S{p{1.3cm}} Sl}%
%{$\,$} & \ensuremath{#1} &\nomenklaturCellHeight\parbox{8cm}{\footnotesize{#2}} \ensuremath{\footnotesize{#3}}%
%\end{tabular}%
%\end{minipage}%
%%}
%\\ \noindent%
%}
%
%%
%%% Start nomenclature below equations - use this command for the first line of each nomenclature-block.
%%
%\newcommand\ehnomenklaturstart[3]{%
%\noindent
%\nomenclature{$#1$}{#2}{}{}%
%\begin{minipage}[top]{0.7\textwidth}\flushleft%\renewcommand\arraystretch{1.25}% (Value=1.0 is for standard spacing}%
%\begin{tabular}{>{\setlength{\parindent}{-0.2cm}}S{p{1.0cm}}S{p{1.3cm}} Sl}
%$\quad$ \\
%where:\nomenklaturCellHeight & \ensuremath{#1} &\nomenklaturCellHeight\parbox{8cm}{\footnotesize{#2}} \ensuremath{\footnotesize{#3}}%
%\end{tabular}%
%\end{minipage}%
%%}
%\\ \noindent%
%}


%
%% Command for definition of new terms in text
%
\newcommand\term[1]{%
    \emph{#1}%
}

%% Command for creating a double-bar over a matrix symbol
\newcommand\dbar[1]{%
    \ensuremath{\bar{\bar{#1}}}%
}

%% Command for definition of planar cross product
\newcommand\planarCross{%
%    \ensuremath{\underline{\times}}%
%    \ensuremath{\dtimes}%
\begin{turn}{90}\ensuremath{\ltimes}\end{turn}
}



    \usepackage{todonotes}


% new commands for individually colored todo 
\newcommand{\todojohnny}[2]{
\todo[#1, color=green!40]{#2 \\- Johnny}
}
\newcommand{\todosimon}[2]{
\todo[#1, color=red!40]{#2 \\- Simon}
}
\newcommand{\todofrank}[2]{
\todo[#1, color=yellow!40]{#2 \\- Frank}
}
\newcommand{\todokim}[2]{
\todo[#1, color=purple!40]{#2 \\- Wex}
}
\newcommand{\todorahlff}[2]{
\todo[#1, color=gray!40]{#2 \\- Rahlff}
}
\newcommand{\todojens}[2]{
\todo[#1, color=blue!40]{#2 \\- Jens Frost}
}
\newcommand{\todomorten}[2]{
\todo[#1, color=orange!40]{#2 \\- Morten}
}

% Color package. Use with listings package and use for color a table
%\usepackage[usenames,dvipsnames,svgnames,table]{xcolor}

    %%%%%%%%%%%%%%%%%%%%%%%%%%%%%%%%%%%%%%%%%%%%%%%%%%%%%%%%%%%%%%%%%%%%%%%%%%%%%
%                            Custom for funktioner                          %
%                     			      	        		                    %
%%%%%%%%%%%%%%%%%%%%%%%%%%%%%%%%%%%%%%%%%%%%%%%%%%%%%%%%%%%%%%%%%%%%%%%%%%%%%


% Package for wrapping figures in text
\usepackage{wrapfig}

% Set the path where LaTeX looks for pictures.
\graphicspath{{../figurer/}}

% You need a newsubfloat element to use subcaption
\newsubfloat{figure}

% Command to set caption styles
\captionnamefont{\bfseries\small}
\captiontitlefont{\itshape\small}
\subcaptionlabelfont{\bfseries\small}
\subcaptionfont{\itshape\small}

% When we don't use chapter numbering, 
% the figure numbering are shown like 
% fx 0.3 if we don't use this line.
\counterwithout{figure}{chapter}


% Command \figur{filename}{caption}{label}{width} 
% for inserting a new figure.
\newcommand{\figur}[4]{ 
 \begin{figure}[ht] 
 \centering 
 \includegraphics[width=#4\textwidth]{#1} 
 \caption{#2} 
 \label{#3} 
 \end{figure}
 } 
 
% Command \wrapfigur{filename}{caption}{label}{width}{r/l}
% for inserting a new figure wraped in the text. 
\newcommand{\wrapfigur}[5]{ 
 \begin{wrapfigure}{#5}{#4 \textwidth}
   \begin{center}
    \includegraphics[width=#4\textwidth]{#1}
   \end{center}
   \caption{#2}
   \label{#3} 
 \end{wrapfigure}
}
 
% Command \subfigur{filename1}{filename2}{caption}
% {subcaption1}{subcaption2}{label1}{label2}
% {width1}{width1} for inserting a new figure.
\newcommand{\subfigur}[9]{
\begin{figure}[ht]
\centering
\subbottom[{#4}\label{#6}]%
    {\includegraphics[width=#8\textwidth]{#1}}\hfill
\subbottom[{#5}\label{#7}]%
    {\includegraphics[width=#9\textwidth]{#2}}
\caption{#3}
\label{fig:subfigures}
\end{figure}
}
% Command \truefigure{filename}{caption}{label}
% for inserting a picture without scaling.
\newcommand{\truefigure}[3]{ 
 \begin{figure}[ht] 
 \centering 
 \includegraphics{#1} 
 \caption{#2} 
 \label{#3} 
 \end{figure}
 } 

% Keeps floats in the section. 
% Help control place to put pix
\usepackage[section]{placeins}

% Commands for rightaligning wraped text in a tabular
\newcommand{\rr}{\raggedright}
\newcommand{\tn}{\tabularnewline}

\usepackage{multirow}

%Tabular pakke som Wex har fundet, men som vi ikke bruger alligvel
%\usepackage{tabu}

% Include code pieces from files
% Documentation:
% ftp://ftp.tex.ac.uk/tex-archive/macros/latex/contrib/listings/listings.pdf
\usepackage{listings}


%Package with enumirate command with extra options (used within cells of tabular)
\usepackage[inline]{enumitem}
    \usepackage{geometry}
\usepackage{amsmath}
\usepackage[some]{background}


\definecolor{titlepagecolor}{cmyk}{1,.60,0,.40}

\DeclareFixedFont{\bigsf}{T1}{phv}{b}{n}{1.5cm}

\backgroundsetup{
scale=1,
angle=0,
opacity=1,
contents={\begin{tikzpicture}[remember picture,overlay]
 \path [fill=ThemeColor] (-0.5\paperwidth,5) rectangle (0.5\paperwidth,10);  
\end{tikzpicture}}
}
\makeatletter                       
\def\printauthor{%                  
    {\large \@author}}              
\makeatother
\author{%
201270762, René R. Jensen \\
\vspace*{0.75cm}
201270763, Sander Louring \\
\vspace*{0.75cm}
201206094, Kasper H. Liholm \\
\vspace*{0.75cm}
201271015, Christian Sørensen \\
\vspace*{0.75cm}
201270893, Andreas H. Jakobsen \\
\vspace*{0.75cm}
201270097, Simon Ø. Kristensen \\
\vspace*{0.75cm}
201270987, Christoffer B. Rasmussen \\



\vspace*{0.75cm}
Vejleder: Steen F. Krøyer
}



    
    
    
    
% % % % % % % % % % % % % % ********% % %Generelt design % % % % % % % % % % % % % % % % % % %
 
%\usepackage[table,dvipsnames]{xcolor} % 'table' er i stedet for \usepackage{colortbl} og 'dvipsnames' er et farveskema hvor man angiver farver med navne
 
\usepackage{xcolor} % farver kan indsætttes i henhold til RGB med definecolor

%Valg af temafarve % Denne farve bruges i chapterstyle %
%\definecolor{ThemeColor}{RGB}{224, 32, 28} % Red
\definecolor{ThemeColor}{RGB}{255,110,0} % Orange

% Valg af font %
\usepackage[sc]{mathpazo}
\linespread{1.05}         % Palatino needs more leading (space between lines)
\usepackage[T1]{fontenc}


% Pakker inkludere overordnede designvalg. Pakkerne er valgt for overskuelighed %
	 %%%%%%%%%%%%%%%%%%%%%%%%%%%%%%%%%%%%%%%%%%%%%%%%%%%%%%%%%%%%%%%%%%%%%%%%%%%%%
%                             Chapter and Section                           %
%                     			      Style	       		                    %
%%%%%%%%%%%%%%%%%%%%%%%%%%%%%%%%%%%%%%%%%%%%%%%%%%%%%%%%%%%%%%%%%%%%%%%%%%%%%


% % % % Redefined so that header will show even on chapter start. % % % %
\makeatletter
    \let\stdchapter\chapter
    \renewcommand*\chapter{%
    \@ifstar{\starchapter}{\@dblarg\nostarchapter}}
    \newcommand*\starchapter[1]{%
        \stdchapter*{#1}
        \thispagestyle{fancy}
        \markboth{\MakeUppercase{#1}}{}
    }
    \def\nostarchapter[#1]#2{%
        \stdchapter[{#1}]{#2}
        \thispagestyle{fancy}
    }
\makeatother




\usepackage{titlesec} % Pakke til opsætning af chapter og sektioner


% % defining new font for the chapternumber. Calling it FuckingHuge 
\newcommand{\FuckingHuge}{% Needs a 'stretchable' font
      \usefont{\encodingdefault}{\rmdefault}{b}{n}%
      \fontsize{35}{45}\normalfont\color{ThemeColor}} %60 80


%%% TAL I OVERSKRIFTER, både chapters og sections. 
    \maxsecnumdepth{subsubsection} %subsubsection er dybdeniveau for numering.
   	\setsecnumdepth{subsubsection}


  
 % % % % CHAPTER OPSÆTNING % % % % % % % % % % % % % % % % % % % % %

%Hang thoose numbers! I margin med tallene for både sections og chapters! En lækker chapter style!
\chapterstyle{hangnum}
 
%Modificering af ovenstående style
%Placering af chapter overskriften på siden i forhold til sidens start %
\renewcommand*{\chapterheadstart}{\vspace*{-20pt}}


% FARVE PÅ CHAPTER % %
%TItel farve og font valg
\renewcommand*{\chaptitlefont}{\normalfont\Huge\color{ThemeColor}}
%Numering farve og font valg
\renewcommand{\chapnumfont}{\normalfont\FuckingHuge\color{ThemeColor}}



% % % % SECTION OPSÆTNING % % % % % % % % % % % % % % % % % % % %
\hangsecnum % vi smidder numering ud i marken

%Vi benytter tema farven til numeringen
\renewcommand\thesection{\color{ThemeColor}\thechapter.\arabic{section}}



    
    


	 \input{Styles/header.tex}
	 \usepackage{listings}


\usepackage[T1]{fontenc}
\usepackage{inconsolata}


\definecolor{bluekeywords}{rgb}{0.13,0.13,1}
\definecolor{greencomments}{rgb}{0,0.5,0}
\definecolor{redstrings}{rgb}{0.9,0,0}

\usepackage{listings}
\lstset{language=[Sharp]C,
  showspaces=false,
  showtabs=false,
  breaklines=true,
  showstringspaces=false,
  breakatwhitespace=true,
  escapeinside={(*@}{@*)},
  commentstyle=\color{greencomments},
  keywordstyle=\color{bluekeywords},
  stringstyle=\color{redstrings},
  %basicstyle=\ttfamily
  basicstyle=\fontsize{10}{13}\selectfont\ttfamily
}

	 

% Setup the margins:
\setlrmarginsandblock{1in}{*}{1.5}
\setulmarginsandblock{3cm}{4cm}{*}



% Setup header and footer: %
% \setheadfoot{headheight}{footskip}
% \setheaderspaces{headdrop}{headsep}{ratio}
\setheadfoot{45pt}{1cm}
\setheaderspaces{1cm}{*}{*}
\checkandfixthelayout
% Se dette pdf for dokmentaton:
% www.tug.dk/_media/foredrag/2003-12-03/slides-src.pdf




 \checkandfixthelayout[nearest]
    \checkthelayout
    \fixthelayout
    \baselineskip \topskip



% % % % % % % % % % % % % % % % % % % %***** % % % % % % % % % % % % % % % % % % % % % % % % %

%BibTex ops�tning
\usepackage[semicolon,authoryear,square]{natbib}

\bibliographystyle{plainnat}
%%%%%%%%%%%%%%%%%%%%%%%%%%%%%%%%%%%%%%%%%%%%%%%%%%%%%%%%%%%%%%%%%%%%%%%%%%%%%%%%%%%%%%%%%%%%%%

   
    
% % Smart listing function. Documentation: http://texblog.org/tag/makenomenclature/
\usepackage{nomencl}
\makenomenclature 



% % % % % % % % % % % % % % % % Reference opsætning % % % % % % % % % % % % % % % % % % % % % % %

\usepackage{varioref} % for references, http://www.ctex.org/documents/packages/bibref/varioref.pdf
   
   
    
% External reference package, zref makes the references into hyperlinks.
\usepackage{url}
\usepackage{hyperref}
%\usepackage{zref-xr}
\usepackage{xr}
%Reference til ekstern Rapport.
\externaldocument[Krav-]{../Kravspecifikation/Kravspecifikation}
\externaldocument[Design-]{../Designdokument/Designdokument}
\externaldocument[Test-]{../Accepttestspecifikation/Accepttestspecifikation}

% % % % % % % % % % % % % % % % % % % % % % % % % % % % % % % % % % % % % % % % % % % % % % % % % % %

%\usepackage{showlabels}


% footruleskip is allready defined in memoir, so we have to undefine it, in order to include fancyhdr for our footer and header.
\let\footruleskip\undefined
\usepackage{fancyhdr}% For making fancy headers http://ctan.org/pkg/fancyhdr



\AtBeginDocument{\renewcommand\contentsname{Indholdsfortegnelse}}

% % % % % % % % % % % % % Inkluderinger af ønskede kapitler % % % % % % % % % % % % %


\begin{document}
\thispagestyle{empty} % No header and footer on this page

\BgThispage
\newgeometry{left=1cm,right=4cm}
\vspace*{2cm}
\noindent
\textcolor{white}{\bigsf Mingle \\ Kravspecifikation}
\vspace*{2.5cm}\par
\noindent
\begin{minipage}{0.45\linewidth}
    \begin{flushright}
        \printauthor
    \end{flushright}
\end{minipage} \hspace{30pt}
%
\begin{minipage}{0.02\linewidth}
    \rule{1pt}{275pt}
\end{minipage} \hspace{-10pt}
%
\begin{minipage}{0.6\linewidth}
\vspace{5pt}
 

\includegraphics[width=11cm]{Logo}
  
Gruppe 2 \newline

Dato 27/5 - 2014 \newline

Ingeniørhøjskolen Aarhus Universitet \newline



\end{minipage}
\restoregeometry
\mbox{} % Insert a blank page so, toc appears on an odd pagenumber
\thispagestyle{empty} % No header and footer on this page
\newpage

\thispagestyle{empty} % No header and footer on this page
\tableofcontents*
\thispagestyle{empty} % No header and footer on this page

% Setup header and footer
\pagestyle{fancy}
\fancyhf{} % Clear all header an footer fields 
\renewcommand{\headrulewidth}{0pt} % Delete headruler
\fancyhead[RO]{\includegraphics[height=40pt]{header}} % Insert picture in right side of header
%\fancyhead[R]{\rightmark} % insert chapter title in right side of header
\fancyfoot[LE,RO]{\thepage} % insert pagenumber on right side of even page, left side of odd page 
\setcounter{page}{1}



\thispagestyle{fancy}
\chapter{Opgavebeskrivelse}
Målet med projektet er i al sin enkelhed, at undersøge mulighederne for at lave EMG signal genkendelse på armens muskler. Hensigten er her at lave genkendelsesalgoritmer ud fr machnelearning principper. f.eks. som deep learning, til at skelne mellem håndbevægelser.\\

Det er i projektet, tanken at anvende et wearable device til opsamling af EMG data. Denne enhed er Myo \href{https://www.thalmic.com/en/myo/}{https://www.thalmic.com/en/myo/}, et EMG armbånd med 8 EMG sensore og 9 aksial IMU med accelaometer, gyroscope og magnetometer. EMG dataen skal bruges til at “træne” programmet/softwaren/alogritmerne til at skelne mellem forskellig håndbevægelse.\\

Projektet er inspireret af myoelektriske proteser som anvender EMG signaler til at styre elektriske armproteser. Det er således håbet at projektet vil kunne anvendes i realtime til at styre armproteser, robotarme, objekter i virtual reality eller lignende.

\begin{itemize}
	\item Anskaffelse af store mængder EMG data fra arm eller lign.
	\item Udvikle visuelt softwareinterface til dataprocessering.
	\item Udvikle machine learning algoritme til genkendelse af EMG data.
	\item Identificere mulige anvendelsesscenarier.
\end{itemize}
\thispagestyle{fancy}
\chapter{Kravspecifikation}
I denne del af bachelor-forbedrelsesjournalen vil projektets foreløbige funktionelle krav blive gennemgået. Kravene, der her bliver gennemgået, skal ses som midlertidige udkast, da projektets mål ikke er endeligt fastlagt.
\\ \\
I projektets context diagram vil systems overordnede struktur blive beskrevet. \\
Foreløbig er intentionen at projektet opdeles i tre særskilte del:\\
\begin{enumerate}
	\item En dataopsamlings del, hvor data kan opsamles, organiseres og glemmes, så det kan anvendes af de øvrige dele af projektet
	\item En machine-lerning-del, hvor den indsamlede data processeres. Videre skal der i denne del være mulighed for at begrænse, hvilke dele af den indsamlede data, der processeres
	\item En form for SDK-del, der kan anvendes i software, hvor der er behov for håndbevægelsesgenkendelse
	
\end{enumerate}
\section{User case diagram}

\section{Fully dressed use cases}
\subsection{UC 1 - Indsaml datasæt}
\begin{table}[h]
	\begin{tabular}{|p{5cm}|p{9cm}|}
		\hline
		\textbf{Navn} & Indsaml dataset \\ \hline
		\textbf{Use case ID} & U1 \\ \hline
		\textbf{initiering} & Admin initierer denne use case \\ \hline
		\textbf{Aktører} & Admin (Primær Aktør ) \\ & Testperson (Sekundær aktør) \\ \hline
		\textbf{samtigdige forekomster} & - \\ \hline
		\textbf{Prækondition} & Testperson har Myo båndet på og dataindsamlings software kører \\ \hline
		\textbf{Postkondition} & Data set er indsamlet \\ \hline
	\end{tabular}
\end{table}
\textbf{Hovedscenarie}
\begin{enumerate}
	\item Admin trykker på knap på at starte for at starte dataindsamlingen.
	\item Testpersonen laver håndtegn.
	\item Admin trykker på knap for at stoppe dataindsamlingen, når håndtegnsbevægelsen er slut.
	\item Admin vælger, hvilken bevægelsestype der er blevet indsamlet og angiver hvem testperson.
	\item Admin trykker på knap for at gemme datasættet.
\end{enumerate}

\subsection{UC 2 - Tilføj ny bevægelsestype}
\begin{table}[h]
	\begin{tabular}{|p{5cm}|p{9cm}|}
		\hline
		\textbf{Navn} & Tilføj ny bevægelsestype \\ \hline
		\textbf{Use case ID} & U2 \\ \hline
		\textbf{initiering} & Admin initierer denne use case \\ \hline
		\textbf{Aktører} & Admin (Primær Aktør ) \\ \hline
		\textbf{samtigdige forekomster} & - \\ \hline
		\textbf{Prækondition} & Dataindsamlings softwaren kører \\ \hline
		\textbf{Postkondition} & En ny bevægelsestype er tilføjet til Dataindsamlings softwaren \\ \hline
	\end{tabular}
\end{table}
\textbf{Hovedscenarie}
\begin{enumerate}
	\item Admin trykker på knap for at tilføje ny bevægelsestype.
	\item Admin angiver title for den nye bevægelsestype.
	\item Admin tildeler evt. farve til den nye bevægelsestype.
	\item Admin vælger, hvilken bevægelsestype der er blevet indsamlet og angiver hvem testperson.
	\item Admin trykker på knap for at gemme bevægelsestypen.
\end{enumerate}

\subsection{UC 3 - Slet datasæt}
\begin{table}[h]
	\begin{tabular}{|p{5cm}|p{9cm}|}
		\hline
		\textbf{Navn} & Slet datasæt \\ \hline
		\textbf{Use case ID} & U3 \\ \hline
		\textbf{initiering} & Admin initierer denne use case \\ \hline
		\textbf{Aktører} & Admin (Primær Aktør ) \\ \hline
		\textbf{samtigdige forekomster} & - \\ \hline
		\textbf{Prækondition} & Et datasæt skal eksistere \\ \hline
		\textbf{Postkondition} & Datasæt er slettet \\ \hline
	\end{tabular}
\end{table}
\textbf{Hovedscenarie}
\begin{enumerate}
	\item Admin vælger eksisterende datasæt.
	\item Admin trykker på knap for at slette valgte datasæt.
\end{enumerate}

\subsection{UC 4 - Verificering af ML genkendelse}
\begin{table}[h]
	\begin{tabular}{|p{5cm}|p{9cm}|}
		\hline
		\textbf{Navn} & Verificering af ML genkendelse \\ \hline
		\textbf{Use case ID} & U4 \\ \hline
		\textbf{initiering} & Admin initierer denne use case \\ \hline
		\textbf{Aktører} & Admin (Primær Aktør ) \\ \hline
		\textbf{samtigdige forekomster} & - \\ \hline
		\textbf{Prækondition} & ML har kategoriseret bevægelse \\ \hline
		\textbf{Postkondition} & Admin har varificeret kategorisering  \\ \hline
	\end{tabular}
\end{table}
\textbf{Hovedscenarie}
\begin{enumerate}
	\item Admin tjekker ML kategoriseringer.
	\begin{enumerate}
		\item Korrekt kategorisering: admin verificerer kategorisering.
		\begin{enumerate}
			\item Verificering gemmes.
		\end{enumerate}
		\item Ukorrekt verificiering: admin ændre kategorisering.
		\begin{enumerate}
			\item Ny kategorisering gemmes som verificeret.
		\end{enumerate}
	\end{enumerate}
\end{enumerate}

\subsection{UC 5 - Anvendelse af bevægelsesgenkendelse}
\begin{table}[h]
	\begin{tabular}{|p{5cm}|p{9cm}|}
		\hline
		\textbf{Navn} & Anvendelse af bevægelsesgenkendelse \\ \hline
		\textbf{Use case ID} & U5 \\ \hline
		\textbf{initiering} & Softwareudvikleren initierer denne use case \\ \hline
		\textbf{Aktører} & Softwareudvikleren (Primær Aktør ) \\ \hline
		\textbf{samtigdige forekomster} & - \\ \hline
		\textbf{Prækondition} & Softwareudvikleren har bevægelsesgenkendelses SDK\\ \hline
		\textbf{Postkondition} & Softwareudvikleren har anvendt bevægelsesgenkendelsen\\ \hline
	\end{tabular}
\end{table}
\textbf{Hovedscenarie}
\begin{enumerate}
	\item Softwareudvikleren implementerer SDK i egen software
	\item Softwareudvikleren anvender SDK til bevægelsesgenkendelse
\end{enumerate}

\subsection{UC 6 - Genkendelse af bevægelse}
\begin{table}[h]
	\begin{tabular}{|p{5cm}|p{9cm}|}
		\hline
		\textbf{Navn} & Genkendelse af bevægelse  \\ \hline
		\textbf{Use case ID} & U6 \\ \hline
		\textbf{initiering} & Testpersonen initierer denne use case \\ \hline
		\textbf{Aktører} & Testperson (Primær Aktør )\\ & Softwareudvikler (Sekundær aktør) \\ \hline
		\textbf{samtigdige forekomster} & - \\ \hline
		\textbf{Prækondition} & Testperson har Myo bånd på, som er forbundet til softwareudviklerens software \\ \hline
		\textbf{Postkondition} & Håndbevægelse er blevet genkendt\\ \hline
	\end{tabular}
\end{table}
\textbf{Hovedscenarie}
\begin{enumerate}
	\item Testperson laver håndbevægelse
	\item Myo-data sendes til softwareudviklerens software på computer el.lign.
	\item Softwareudviklerens software processer Myo-data
	\item Myo-data genkendes
	\item Testpersonens håndbevægelse genkendes
\end{enumerate}
\thispagestyle{fancy}
\chapter{Projektplan}
\subsection*{Evt. eksperimenter}
Herunder er en liste af undersøgelser, der kan være interessante at lave i forbindelse med projektet.
\begin{itemize}
	\item Hvor meget kalibrering er nødvendigt?
	\begin{itemize}
		\item Undersøgelse af hvor meget kalibrering der er nødvendigt for at få brugbare dataset fra armbåndet.
	\end{itemize}
	\item Hvor mange sensorer er nødvendige?
	\begin{itemize}
		\item Undersøgelse af hvor få sensorer man kan nøjes med at bruge og stadig få brugbare dataset fra armbåndet.
	\end{itemize}
	\item Kan bevægelser genkendes på tværs af personer?
	\begin{itemize}
		\item Undersøgelse af om den indsamlede data er generisk på tværs af personer. Vil den enkelte bevægelse kunne genkendes fra én person til en anden.
	\end{itemize}
	\item Kan personer genkendes
	\begin{itemize}
		\item Undersøgelse af data indsamlet fra flere personer kan skelne mellem hver person ud fra deres bevægelse.
		\end{itemize}
\end{itemize}

\subsection*{Teknologier}
Herunder findes en liste over relevante teknologier el.lign. der muligvis kan anvendes i projektet
\begin{itemize}
	\item Artificial neural network
	\item Google Cloud Computing
	\item Microsoft Azure - Machine Learning
	\item Parse.com
	\item Unity - VR
\end{itemize}

\subsection*{Relevante kilder}
Herunder er en liste med interessante online kurser. Videre forventes det, at yderligere materiale, som bøger og artikler kan findes gennem AU Library, \href{http://library.au.dk/materialer/}{http://library.au.dk/materialer/}
\begin{itemize}
	\item E-kursus i Artificial intelligence
		\begin{itemize}
			\item \href{https://www.udacity.com/course/intro-to-artificial-intelligence--cs271}{https://www.udacity.com/course/intro-to-artificial-intelligence--cs271}
			\item \href{https://www.udacity.com/course/knowledge-based-ai-cognitive-systems--ud409}{https://www.udacity.com/course/knowledge-based-ai-cognitive-systems--ud409}
		\end{itemize}
	\item E-kursus i machine learning
		\begin{itemize}
			\item \href{https://www.udacity.com/course/intro-to-machine-learning--ud120}{https://www.udacity.com/course/intro-to-machine-learning--ud120}
			\item \href{https://www.udacity.com/course/machine-learning--ud262}{https://www.udacity.com/course/machine-learning--ud262}
		\end{itemize}

\end{itemize}



\thispagestyle{fancy}
\chapter{Konklusion}
\label{chp:konklusion}
Yo!


%\bibliographystyle{plainnat}
\bibliography{bibliografi}% Selects .bib file AND prints bibliography
\thispagestyle{fancy}

%\listoftodos % Make a list of todo's
\end{document}